\part{Phonology}

\chapter{Phonemic inventory}

\section{Vowels}

The Interturkic vowel system is characterised by oppositions in backness, rounding, and height.
These oppositions are vitally important to the process of vowel harmony, and the arrangement of Interturkic\apostrophe{}s nine vowel phonemes along said oppositions in shown below.

\begin{smalltable}{Vowel phonemes}{c|c|c|c|c}
	& \multicolumn{2}{c|}{\textbf{Front}} & \multicolumn{2}{c|}{\textbf{Back}} \\
	\cline{2-5}
	& \textbf{unrounded} & \textbf{rounded} & \textbf{unrounded} & \textbf{rounded} \\
	\hline
	\textbf{High} & i & y & ɯ & u \\
	\hline
	\textbf{Low} & ɛ æ & œ & ɑ & ɔ \\
\end{smalltable}

Dipthongs are not considered to be separate phonemes.
Sequences perceived as rising dipthongs are analysed as /j/ or /w/ followed by a vowel (see example~\ref{ex:rising-diphthongs}).
Sequences perceived as falling dipthongs are analysed as a vowel followed by /j/ or /w/ (see example~\ref{ex:falling-diphthongs}).

\begin{lingexample}
	\ex \label{ex:rising-diphthongs} /jɯl/ `year' \\
		/wætæn/ `homeland'
\end{lingexample}

\begin{lingexample}
	\ex \label{ex:falling-diphthongs} /ɑj/ `moon, month' \\
		/birɛw/ `someone, somebody'
\end{lingexample}

\section{Consonants}

Interturkic has 26 consonant phonemes, primarily classified by their voicing.
The unvoiced consonants are /p, t, k, q, ʔ, t͡ʃ, f, s, ʃ, x, h/.
The voiced consonants are /m, n, ŋ, b, d, g, d͡ʒ, v, z, ʒ, ɣ, w, l, j, ɾ/.

\begin{smalltable}{Consonant phonemes}{c|c|c|c|c|c}
	& \textbf{Labial} & \textbf{Alveolar} & \textbf{Postalveolar} & \textbf{Dorsal} & \textbf{Glottal} \\
	\hline
	\textbf{Nasal} & m & n & & ŋ & \\
	\hline
	\textbf{Stop} & p b & t d & & k q g & ʔ \\
	\hline
	\textbf{Affricate} & & & t͡ʃ d͡ʒ & & \\
	\hline
	\textbf{Fricative} & f v & s z & ʃ ʒ & x ɣ & h \\
	\hline
	\textbf{Approximant} & w & l & & j & \\
	\hline
	\textbf{Tap} & & ɾ & & & \\
\end{smalltable}

\chapter{Phonotactics}

\chapter{Morphophonological alternations}

\section{Final consonant lenition}

\section{Vowel harmony}

\section{Consonant assimilation}

\chapter{Stress}
